\subsection{子代数。理想。商代数}

设 A 是一个交换环，E 是一个 A-代数。若 F 是 E 的一个子 A-模，且在 E 的乘法下稳定，则 E 的乘法在 \(F \times F\) 上的限制（连同 F 上的 A-模结构）在 F 上定义了一个 A-代数结构。带有此结构的 F 称为 A-代数 E 的一个子代数。E 的子代数的任一交都是 E 的一个子代数。对 E 的元素的每个族 \((x_{i})_{i \in I}\) ，包含所有 \(x_{i}\) 的 E 的子代数之交称为由族 \((x_{i})_{i \in I}\) 生成的 E 的子代数，而 \((x_{i})_{i \in I}\) 称为该子代数的一个生成系（或生成族）。若 \(u: E \to E'\) 是一个 A-代数同态，则 E 的每个子代数 F 的像 \(u(F)\) 都是 \(E'\) 的子代数。

设E是一个结合代数。对于E的每个子集M，E中与M的所有元素都可交换的元素所构成的集合M'是E的一个子代数，称为M在E中的中心化子代数（I，§1，第5号）。M'在E中的中心化子M''也称为M的双中心化子；显然M ⊆ M''。由此可知M'包含于它的双中心化子M''中，而M''正是M''的中心化子；但关系M ⊆ M''蕴含M'' ⊆ M'，所以M' = M''（参见《集合论》，III，§1，第7号，命题2）。若F是E的子代数，则F的中心是F与其在E中的中心化子F'的交F ∩ F'。注意若

F 是交换的，则 \(F \subset F'\) ，从而 \(F' \supset F''\) ；此时 F 的双中心化子 \(F''\) 就是 F 的中心。

对于某些非结合代数（例如李代数），子代数的中心化子与中心这两个概念另有定义（《李群与李代数》，I，§ 1，第 6 号）。

子集 \(\mathfrak{a}\) ——A-代数 \(\mathbf{E}\) 的子集——称为 \(\mathbf{E}\) 的左理想（相应地，右理想），如果 \(\mathfrak{a}\) 是 \(\mathbf{E}\) 的子 A-模，且关系 \(x \in \mathfrak{a}, y \in \mathbf{E}\) 蕴含 \(yx \in \mathfrak{a}\) （相应地， \(xy \in \mathfrak{a}\) ）。这等同于说 \(\mathfrak{a}\) 是 \(\mathbf{E}\) 的左理想或反代数 \(\mathbf{E}^0\) 的右理想。\(\mathbf{E}\) 的双边理想是指一个子集 \(\mathfrak{a}\) ，它既是 \(\mathbf{E}\) 的左理想又是右理想。当 \(\mathbf{E}\) 结合且具有单位元 \(e\) 时，则对 \(\alpha \in \mathbf{A}\) 与 \(x \in \mathbf{E}\) ，由（2）（第 1 号）得 \(\alpha x = (\alpha e)x = x(\alpha e)\) ，因而环 \(\mathbf{E}\) （I，§ 8，第 6 号）的（右、左、双边）理想与代数 \(\mathbf{E}\) 的（右、左、双边）理想相同。代数 \(\mathbf{E}\) 的左（相应地，右、双边）理想的每个和与每个交都是左（相应地，右、双边）理想。包含子集 \(\mathbf{X}\) （\(\mathbf{E}\) 中的）的左（相应地，右，相应地，双边）理想之交称为 \(\mathbf{E}\) 的由 \(\mathbf{X}\) 生成的左（相应地，右，相应地，双边）理想。

设 \(\mathfrak{b}\) 是 A-代数 E 的一个双边理想。若 \(x \equiv x'\pmod{\mathfrak{b}}\) 和 \(y \equiv y'\pmod{\mathfrak{b}}\) ，则

\[
x (y - y ^ {\prime}) \in \mathfrak {b} \quad \text { and } \quad (x - x ^ {\prime}) y ^ {\prime} \in \mathfrak {b}
\]

则 \(xy \equiv x'y' \pmod{b}\) 。因而可在商 A-模 E/b 上定义一个内法则，它是 E 的乘法法则 \((x,y) \mapsto xy\) 关于等价关系 \(x \equiv x' \pmod{b}\) 的商法则（I，§ 1，第 6 号）。立即可以验证，这个商法则是 \((\mathrm{E} / \mathfrak{b}) \times (\mathrm{E} / \mathfrak{b})\) 到 E/b 中的一个 A-双线性映射；因此，它连同 E/b 上的 A-模结构在 E/b 上定义了一个 A-代数结构。带有这个代数结构的 E/b 称为代数 E 关于双边理想 b 的商代数。典范映射 \(p: \mathrm{E} \to \mathrm{E} / \mathfrak{b}\) 是一个代数同态。

设 E, E' 为两个 A-代数，且 \(u:E \rightarrow E'\) 为一个代数同态。像 \(u(\mathbf{E})\) 是 \(E'\) 的一个子代数，而核 \(b = u(0)\) 是 E 的一个双边理想；进一步，在 u 的典范分解中：

\[
\mathrm{E} \xrightarrow {p} \mathrm{E} / \mathfrak {b} \xrightarrow {v} u (\mathrm{E}) \xrightarrow {j} \mathrm{E} ^ {\prime}
\]

v 是一个代数同构。更一般地，当把“环”一词换成“代数”时，第 I 章 § 8 第 9 号的所有结果（连同它们的证明）仍然成立。

设 A 是一个交换环，E 是一个 A-代数。在集合 \(\tilde{\mathbf{E}} = \mathbf{A} \times \mathbf{E}\) 上，我们定义以下合成法则：

\[
\begin{array}{r} (\lambda , x) + (\mu , y) = (\lambda + \mu , x + y) \\ (\lambda , x) (\mu , y) = (\lambda \mu , x y + \mu x + \lambda y) \\ \lambda (\mu , x) = (\lambda \mu , \lambda x). \end{array}
\]

立即可验证， \(\tilde{E}\) 带这些合成律是 A 上的一个代数，且 \((1,0)\) 是这个代数的一个单位元。集合 \(\{0\}\times E\) 是 \(\tilde{E}\) 的一个双边理想，且 \(x\mapsto(0,x)\) 是代数 E 到子代数 \(\{0\}\times E\) 上的一个同构，借助它，E 与 \(\{0\}\times E\) 等同。 \(\tilde{E}\) 称为由 E 添加单位元而得到的代数；它为结合的（相应地，交换的）当且仅当 E 为结合的（相应地，交换的）。

